\emph{Setting.} The script checks that the Fock module carries the Clifford relations of the ladder pairs, with $v_{j}$ the wedge, $w_{j}$ the contraction and $\dagger$ the adjoint for the inner product in which the occupation basis is orthonormal (Theorem \claimref{tower.fock-module-one-qubit-per-rung}, Proposition \claimref{tower.ladder-pairs-satisfy-car-vectors}); a vector state is linear, normalised, real and positive, so it satisfies Premise \claimref{synthesis.premise-positivity-of-values}. With $h_{j} = v_{a_{j}}w_{b_{j}} - v_{b_{j}}w_{a_{j}}$: $h_{j}^{\dagger} = -h_{j}$ and $h_{j}^{2} = -\Pi_{j}$, so $R_{j}R_{j}^{\dagger} = (1 - \Pi_{j}) + (c^{2} + s^{2})\Pi_{j} = 1$; $R_{j}$ is even and involves only pair $j$, so it commutes with every element that involves no generator of pair $j$; $U_{j}U_{j}^{\dagger} = \mathrm{off}_{0} + \mathrm{on}_{0}R_{j}R_{j}^{\dagger} = 1$. Since $\mathrm{on}_{0}B = 0$, $U_{j}B = B$; since $\mathrm{on}_{0}$ fixes a vector with the system quantum on rung $0$, $A_{K} = \prod_{j \in K}R_{j}\,A_{\emptyset}$.

\emph{(i).} Let $x$ involve no generator of the pairs outside $K$. Then $A_{N} = \bigl(\prod_{j \notin K}R_{j}\bigr)A_{K}$ and $x$ commutes with each $R_{j}$, $j \notin K$, so
\[
\langle A_{N}, xB\rangle = \Bigl\langle A_{K},\, x\prod_{j \notin K}R_{j}^{\dagger}B\Bigr\rangle .
\]
In $B$ pair $j$ holds one quantum, on rung $a_{j}$, so $\Pi_{j}B = B$ and $R_{j}^{\dagger}B = cB - s\,h_{j}B$, where $h_{j}B$ has that quantum on rung $b_{j}$ instead. Expanding the product, every term other than $c^{N-k}B$ has, for at least one $j \notin K$, the quantum of pair $j$ on rung $b_{j}$. In $A_{K}$ that pair is unrotated and has its quantum on rung $a_{j}$, and $x$ does not change the occupation of the rungs of pair $j$; vectors with different occupations of a rung are orthogonal. So $\langle A_{N}, xB\rangle = c^{N-k}\langle A_{K}, xB\rangle$, and the same argument applied to $x^{\dagger}$ gives $\langle B, xA_{N}\rangle = c^{N-k}\langle B, xA_{K}\rangle$. No assumption on the parity of $x$ is used. The two branch vectors each hold $1 + N$ quanta, so an odd $x$ has value zero between them: for odd elements the identity is $0 = 0$. The script verifies the identity for $N = 3$, for $c = \tfrac35$ and $c = 0$, for every proper subset $K$ and every element of a linear basis of the algebra on the system and the pairs of $K$ (products of $1, v, w, vw$ over the rungs, $13072$ elements in all, of which $74$ and $14$ have nonzero $I_{K}$); by linearity this covers every element at that size. For $K$ all $N$ pairs the identity is trivial.

\emph{(ii).} $O_{0}$ is even and self-adjoint, $O_{0}^{2} = \mathrm{on}_{0}\mathrm{off}_{1} + \mathrm{off}_{0}\mathrm{on}_{1}$ is an idempotent, and $U$ is an even unitary; so $O_{N}$ is even and self-adjoint with $O_{N}^{2} = UO_{0}^{2}U^{\dagger}$ idempotent. $\langle\Psi_{N}, O_{N}\Psi_{N}\rangle = \langle\Psi_{0}, O_{0}\Psi_{0}\rangle = 2\cdot\tfrac35\cdot\tfrac45 = \tfrac{24}{25}$. $U^{\dagger}A_{N} = A_{\emptyset}$ and $U^{\dagger}B = B$, and $O_{0}$ has value $0$ in each of $A_{\emptyset}$, $B$ (it moves the system's quantum). The script confirms these for $N = 1, 2, 3$ and both values of $c$, including evenness and the idempotent on the whole module.

\emph{(iii).} Every element of the algebraic limit involves finitely many generators (Theorem \claimref{tower.algebraic-limit-even-charge}), hence the system and a finite set $K$ of pairs; enlarge $K$ if necessary. For $N$ beyond the largest pair in $K$, $\omega_{\Psi_{N}}(x) = \tfrac{9}{25}\langle A_{N}, xA_{N}\rangle + \tfrac{16}{25}\langle B, xB\rangle + I_{N}(x)$; $\langle A_{N}, xA_{N}\rangle = \langle A_{K}, xA_{K}\rangle$ because the $R_{j}$, $j \notin K$, are unitary and commute with $x$; and $I_{N}(x) = c^{N-k}I_{K}(x) \to 0$ since $|c| < 1$, with $I_{N}(x) = 0$ at once when $c = 0$ and at least one coupled pair lies outside $K$. The formula for $\omega_{\Psi_{N}}(x)$ just used needs $N$ beyond the largest pair in $K$, because $\Psi_{N}$ fills only the pairs $j < N$ while $A_{K}$ and $B$ in the definition of $\omega_{A}$, $\omega_{B}$ have the pairs of $K$ filled. So for $c = 0$ equality with the mixture holds when every pair $x$ involves is among the $N$ coupled pairs and at least one of them is not involved by $x$, in particular for $N > p + 1$ with $p$ the largest index $x$ involves. Both conditions matter at $N = 1$: $x = O_{1}$ has the value $\tfrac{24}{25}$ against $0$; $x = \mathrm{on}_{a_{1}}$ has the value $0$ in $\Psi_{1}$, where pair $1$ is empty, against $\tfrac{16}{25}$ in the mixture, where pair $1$ has its quantum on rung $a_{1}$ in branch $B$ and on rung $b_{1}$ in branch $A$. The script computes both. Here $I_{K}(x)$ is computed with the pairs outside $K$ present and unrotated in both branches, and it does not depend on how many such pairs there are, since $x$ does not involve them and they are in the same unit vector in both branches. The value $\langle A_{K}, xA_{K}\rangle$ does not depend on the choice of $K$ by the same commutation. A pointwise limit of states under (P) is linear, normalised, real and positive.

\emph{(iv).} In pair $j$ the quantum is in the mode $(1, 0)$ in branch $B$ and in the mode $(c, -s)$ in branch $A$ (coordinates on the rungs $a_{j}$, $b_{j}$; the script computes the overlap $c$ of the two record vectors, and computes the one-particle operators $\langle v_{i}w_{k}\rangle$ of both branch vectors for $N = 1, 2, 3$: both are projections and the squared Hilbert-Schmidt norm of their difference is $2 + 2s^{2}N$, the $2$ from the system's two rungs). The projections are $\mathrm{diag}(1, 0)$ and the matrix with rows $(c^{2}, -cs)$, $(-cs, s^{2})$; the squared Hilbert-Schmidt norm of their difference is $(c^{2} - 1)^{2} + 2c^{2}s^{2} + s^{4} = 2s^{2}$, positive and the same for every pair, so the sum over infinitely many pairs diverges. Each branch state is the state with a definite set of occupied one-particle modes, a gauge-invariant quasi-free state whose operator is a projection, pure by clause (d) of Fact \claimref{synthesis.imported-powers-stormer-takesaki-disjointness}; for projections the square roots in clause (a) are the projections themselves, so the two states are not quasi-equivalent, and being factor states they are disjoint. This paragraph applies the import to the completed complex algebra and is not proved here.
